% =============================================================================
% Primer Template (Conceptual, Non-Problem Content)
% =============================================================================
% Usage:
% - Copy into: latex/sections/<domain>/primer.tex (or a more specific name)
% - Keep code in external files under: problems/<domain>/snippets/<slug>/
% - From primer.tex, reference code with paths like:
%     ../../../problems/<domain>/snippets/<slug>/solution.cpp
%     ../../../problems/<domain>/snippets/<slug>/solution.py
% - Keep this focused on recognition, invariants, core ops, complexity, pitfalls
.
% =============================================================================

% Replace <Topic> with the data structure / concept name
\subsection{Primer — Arrays & Sequences}

\paragraph{What It Is}
% Concisely describe the structure/concept, its representation, and when it is u
sed.
% Emphasise the mental model (e.g., nodes + pointers; contiguous array + indices
).
An array is a $contiguous$ block of memory storing elements of the same type, supporting $O(1)$ access by index. A dynamic array (e.g. C++ \texttt{std::vector}, Python \texttt{list}) grows by allocating a larger block and copying elements when capacity is exceeded, providing $amortized$ $O(1)$ append. A $sequence$ generalises the idea to any indexable or iterable ordered collection; here we emphasise contiguous arrays.

\paragraph{Recognition Patterns}
% How to identify that this structure is implied by a question. Provide 3–6 cues
.
\begin{itemize}[nosep]
    \item Index math across contiguous elements (prefix/suffix, window, ranges).
    \item Constraints emphasising in-place mutation, $O(1)$ extra space, single pass
    \item Repeated accumulation or difference of segments (prefix sums/scans).
    \item Local adjacency decisions: next/previous comparison, runs, preaks/valleys
\end{itemize}

\paragraph{Core Operations}
% Briefly define the main operations and their expected complexity.
% Include short pseudocode blocks when clarity benefits.
\begin{CodeBlock}[style=pseudo,caption={<Topic>: Core Operations (pseudocode)}]{
18}
access(A, i):
    return A[i] // O(1)

traverse(A):
  for i from 0 to n-1:
    visit(A[i])

insert_at(A, i, x):
  grow A by 1
  for j from n-1 downto i:
    A[j+1] = A[j]
  A[i] = x

erase_at(A, i):
  for j from i to n-2:
    A[j] = A[j+1]
  shrink A by 1

prefix_sums(A):
  P[0] = 0
  for i from 0 to n-1:
    P[i+1] = P[i] + A[i]
  return P

two_pointers_sorted(A, target): 
  i = 0; j = n - 1
  while i < j:
    s = A[i] + A[j]
    if s == target: return (i, j)
    if s < target: i = i + 1 else: j = j - 1
  return not found

sliding_window(A):             
  l = 0
  for r from 0 to n-1:
    add(A[r])         
    while window_invalid():
      remove(A[l])
      l = l + 1
    update_answer()
\end{CodeBlock}

\paragraph{Invariants \& Correctness}
% State the key invariants maintained (e.g., head/tail consistency, ordering, sh
ape).
% Sketch why the core operations preserve these invariants.
Array scans typically maintain index-based invariants:
(i) after $i$ steps, positions $[0..i)$ have been processed,
(ii) two-pointer: indices satisfy $0 \le i \le j < n$ and a monotone move (left increases or right decreases) prevents rework,
(iii) sliding window: the maintained sub-array $[l,r)$ satisfies a property (e.g. distinct elements, sum $\le K$; expansion and contraction preserve the property.

\paragraph{Complexity}
% Summarise asymptotic time/space for each core operation.
\begin{itemize}[nosep]
  \item Access: time $O(1)$; Traversal: $O(n)$
  \item Insert/delete at end (dynamic array): amortised $O(1)$, at index $O(n)$
  \item Prefix sums / scan: $0(n)$ time, $O(n)$ extra space (or $0(1)$ in place if possible)
  \item Two-pointers / sliding window: typically $O(n)$ if each pointer advances monotonically.
\end{itemize}

\paragraph{Common Pitfalls}
% Practical footguns and frequent mistakes.
\begin{itemize}[nosep]
  \item Off-by-one at boundaries; confusing inclusive/exclusive ranges.
  \item Forgetting to update all window state on both expand and shrink
  \item Overflow in sums/producs; use wider types for accumulation.
  \item Accidental $O(n^2)$ from repeated slicing/copying in inner loops.
  \item Assuming sortedness/uniqueness without verifying constraints.
\end{itemize}

% -----------------------------------------------------------------------------
% Example Implementations (external files; adjust paths for your primer location
)
% -----------------------------------------------------------------------------
% From latex/sections/<domain>/primer.tex, use ../../../ to reach repo root.

% \bigskip
% \noindent\textbf{C++ Example}
% \lstinputlisting[
%   style=cpp,
%   caption={<Topic> — Minimal C++ Example}
% ]{../../../problems/<domain>/snippets/<slug>/solution.cpp}

% \bigskip
% \noindent\textbf{Python Example}
% \lstinputlisting[
%   style=python,
%   caption={<Topic> — Minimal Python Example}
% ]{../../../problems/<domain>/snippets/<slug>/solution.py}
